\subsection{折叠标量通道上的 \texorpdfstring{$\pi$}{pi} 精度律、偶 \texorpdfstring{$\zeta$}{zeta} 首差定位与外置预算壁垒}\label{subsec:conclusion-foldgauge-pi-evenzeta-budget-comparison}
经典 \(\pi\) 计算可分解为几何耗尽、AGM、Ramanujan--Chudnovsky、BBP 与低差异采样等若干机制；它们的差异通常表现为不同的辅助状态、模性输入或位地址结构。本文内部的 fold-gauge 接口则走向另一极端：由推论 \ref{cor:fold-bin-recover-2pi}、定理 \ref{thm:fold-bin-gauge-constant-stirling-bernoulli-hierarchy}、定理 \ref{thm:xi-time-part60ab-gauge-constant-recursive-bernoulli-stripping} 与定理 \ref{thm:fold-oracle-capacity-closed-form} 可见，\(2\pi\)、全部偶 Bernoulli 数与全部偶值 \(\zeta(2r)\) 已被压缩进单一标量序列 \(C_m\) 及其预算截断 \(\mathcal C_m(B)\) 之中。沿这条内部接口，可进一步抽取出下列比较型结论：一旦忘却全部辅助状态而只保留折叠标量，任何 \(\pi\)-近似都回落到同一条黄金线性精度律；两条投影兼容序列之间的首个可观测分歧则与偶 \(\zeta\)-塔的首差阶严格对位；而任何试图把该通道提升为远位局域提取机制的方案，都必须支付线性规模的外置地址预算。

\begin{theorem}[折叠标量通道的黄金线性精度律]\label{thm:conclusion-foldgauge-pi-golden-linear-law}
\leanverified{paper\_conclusion\_foldgauge\_pi\_golden\_linear\_law}
记
\[
q:=\frac{\varphi}{2}\in(0,1),
\qquad
\pi_m^{\mathrm{fold}}:=\frac{C_m}{2}.
\]
则存在显式常数
\[
a_1=\frac{1}{3\varphi}\neq 0
\]
使得
\[
\log C_m-\log(2\pi)=a_1 q^m+O(q^{3m}),
\]
\[
\pi_m^{\mathrm{fold}}-\pi=\frac{\pi}{3\varphi}\,q^m+O(q^{2m}).
\]
因此其有效二进精度满足
\[
-\log_2\bigl|\pi_m^{\mathrm{fold}}-\pi\bigr|
=
m\log_2\!\Bigl(\frac{2}{\varphi}\Bigr)+O(1).
\]
等价地，若要求 \(n\) 位二进有效精度，则必要且充分的窗口尺度为
\[
m=
\frac{n}{\log_2(2/\varphi)}+O(1)
\approx 3.27056\,n.
\]
\end{theorem}
\begin{proof}
由推论 \ref{cor:xi-time-part60ab-gauge-constant-leading-mode-ratio}，
\[
\log C_m-\log(2\pi)=a_1 q^m+O(q^{3m}),
\qquad
a_1=\frac{1}{3\varphi}.
\]
记
\[
E_m:=\log C_m-\log(2\pi).
\]
则 \(E_m=a_1 q^m+O(q^{3m})\)，从而
\[
C_m
=
2\pi e^{E_m}
=
2\pi\bigl(1+E_m+O(E_m^2)\bigr)
=
2\pi\bigl(1+a_1 q^m+O(q^{2m})\bigr).
\]
于是
\[
\pi_m^{\mathrm{fold}}-\pi
=
\frac{C_m-2\pi}{2}
=
\pi a_1 q^m+O(q^{2m})
=
\frac{\pi}{3\varphi}\,q^m+O(q^{2m}).
\]
由于前导常数 \(\pi/(3\varphi)\) 非零，故
\[
\bigl|\pi_m^{\mathrm{fold}}-\pi\bigr|
=
\frac{\pi}{3\varphi}\,q^m\bigl(1+O(q^m)\bigr),
\]
取二进对数即得
\[
-\log_2\bigl|\pi_m^{\mathrm{fold}}-\pi\bigr|
=
m\log_2(q^{-1})+O(1)
=
m\log_2\!\Bigl(\frac{2}{\varphi}\Bigr)+O(1).
\]
最后对 \(m\) 反解即得窗口尺度公式。证毕。
\end{proof}

\begin{proposition}[经典加速机制在折叠标量投影下的坍缩]\label{prop:conclusion-foldgauge-classical-pi-collapse}
\leanpartial{paper\_conclusion\_foldgauge\_classical\_pi\_collapse}{给定 Lean 签名量化任意命题且没有对应假设，因此在一致内核中不可证明。}
设 \(\mathcal A\) 为任一用于计算 \(\pi\) 的过程，其在分辨率 \(m\) 上除标量输出外还携带若干辅助状态。定义折叠标量投影 \(P\) 为忘却全部辅助状态而仅保留本文内部标量 \(C_m\) 的操作，即
\[
P(\mathcal A)_m:=C_m.
\]
则投影后的标量 \(\pi\)-近似
\[
\frac{P(\mathcal A)_m}{2}
=
\pi_m^{\mathrm{fold}}
\]
必服从定理 \ref{thm:conclusion-foldgauge-pi-golden-linear-law} 的黄金线性精度律。特别地，AGM 的二次收敛、Ramanujan--Chudnovsky 级数的高位增益，以及 BBP 的远位局域性，都不是该投影的不变量；它们只能存在于被该投影忘却的辅助模态中。
\end{proposition}
\begin{proof}
投影 \(P\) 的定义已强制
\[
\frac{P(\mathcal A)_m}{2}=\frac{C_m}{2}=\pi_m^{\mathrm{fold}}.
\]
因此其误差渐近由 \(C_m\) 的内部结构唯一决定，而定理 \ref{thm:conclusion-foldgauge-pi-golden-linear-law} 已完全固定这一结构。于是任何外部快速机制在纯标量投影下都回落到同一黄金线性误差谱。证毕。
\end{proof}

\begin{theorem}[\texorpdfstring{\(\pi\)}{pi}--偶 \texorpdfstring{\(\zeta\)}{zeta} 塔的首差阶定位]\label{thm:conclusion-foldgauge-pi-evenzeta-first-difference}
\leanverified{paper\_conclusion\_foldgauge\_pi\_evenzeta\_first\_difference}
仍记 \(q=\varphi/2\)。设 \((C_m)_{m\ge 1}\) 与 \((\widetilde C_m)_{m\ge 1}\) 为两组投影兼容（project-compatible）的正序列，即
\[
C_m\to 2\pi,
\qquad
\widetilde C_m\to 2\pi,
\]
并且对每个固定整数 \(R\ge 1\) 都有展开
\[
\log C_m-\log(2\pi)
=
\sum_{j=1}^{R}a_j q^{(2j-1)m}
+O\!\bigl(q^{(2R+1)m}\bigr),
\]
\[
\log \widetilde C_m-\log(2\pi)
=
\sum_{j=1}^{R}\widetilde a_j q^{(2j-1)m}
+O\!\bigl(q^{(2R+1)m}\bigr),
\]
其中 \((a_j)_{j\ge 1}\)、\((\widetilde a_j)_{j\ge 1}\) 由命题 \ref{prop:fold-bin-gauge-bernoulli-extraction-operator} 同型的抽取算子读出。若存在最小整数 \(r\ge 1\) 使
\[
a_j=\widetilde a_j\qquad (1\le j<r),
\qquad
a_r\neq \widetilde a_r,
\]
则有
\[
\log C_m-\log \widetilde C_m
=
(a_r-\widetilde a_r)\,q^{(2r-1)m}
+O\!\bigl(q^{(2r+1)m}\bigr).
\]
特别地，若两序列经同一抽取器恢复出的前 \(r-1\) 个偶值数据一致，而第 \(r\) 个偶值数据首次不同，则首个可观测分歧恰发生在尺度 \(q^{(2r-1)m}\)。
\end{theorem}
\begin{proof}
对给定的 \(r\)，把两条展开都取到 \(R=r\) 阶并逐项相减，得到
\[
\log C_m-\log \widetilde C_m
=
\sum_{j=1}^{r}(a_j-\widetilde a_j)q^{(2j-1)m}
+O\!\bigl(q^{(2r+1)m}\bigr).
\]
由假设 \(a_j=\widetilde a_j\) 对 \(1\le j<r\) 成立，遂化为
\[
\log C_m-\log \widetilde C_m
=
(a_r-\widetilde a_r)q^{(2r-1)m}
+O\!\bigl(q^{(2r+1)m}\bigr),
\]
这就是所求结论。

最后，由命题 \ref{prop:fold-bin-gauge-bernoulli-extraction-operator} 与 Bernoulli--\(\zeta\) 恒等式，抽取系数与偶 Bernoulli 数、偶值 \(\zeta(2r)\) 一一对应；因此“首差系数在第 \(r\) 级出现”与“首差偶值数据在 \(2r\) 级出现”等价。证毕。
\end{proof}

\begin{theorem}[远位局域提取的线性外置地址预算壁垒]\label{thm:conclusion-foldgauge-local-extraction-budget-barrier}
设某个基于折叠标量通道的局域提取方案在尺度 \(m\) 上允许使用至多 \(B_m\) 个外置地址位，并记其成功率为 \(p_m\)。则有上界
\[
p_m\le
\mathrm{Succ}_m(B_m)
:=
\frac{1}{2^m}\sum_{x\in X_m}\min\bigl(d_m(x),2^{B_m}\bigr)
\le
\frac{2^{B_m}}{\bar d_m},
\]
其中
\[
\bar d_m:=\frac{2^m}{|X_m|}
=
\frac{2^m}{F_{m+2}}.
\]
若要求
\[
p_m\longrightarrow 1,
\]
则必有
\[
B_m\ge m\log_2\!\Bigl(\frac{2}{\varphi}\Bigr)+O(1).
\]
特别地，纯粹标量折叠通道 \(B_m\equiv 0\) 满足 \(p_m\to 0\)；因而在不外加地址寄存器的情形下，不可能实现 BBP 型的远位局域提取。换言之，任何想把该通道提升为逐位可寻址机制的方案，都必须支付线性规模的外置地址预算。
\end{theorem}
\begin{proof}
由定理 \ref{thm:fold-oracle-capacity-closed-form}，任何至多使用 \(B_m\) 个外置位的局域提取方案，其成功率都不可能超过最优预言机成功率
\[
\mathrm{Succ}_m(B_m)
=
\frac{1}{2^m}\sum_{x\in X_m}\min\bigl(d_m(x),2^{B_m}\bigr).
\]
又由于
\[
\min\bigl(d_m(x),2^{B_m}\bigr)\le 2^{B_m}
\qquad (\forall x\in X_m),
\]
故
\[
\mathrm{Succ}_m(B_m)
\le
\frac{|X_m|\,2^{B_m}}{2^m}
=
\frac{2^{B_m}}{\bar d_m}.
\]

若 \(p_m\to 1\)，则特别有 \(\mathrm{Succ}_m(B_m)\to 1\)，从而上式强制
\[
2^{B_m}\ge (1-o(1))\,\bar d_m.
\]
而由 Binet 公式，
\[
\bar d_m
=
\frac{2^m}{F_{m+2}}
=
\frac{\sqrt5}{\varphi^2}
\Bigl(\frac{2}{\varphi}\Bigr)^m
\bigl(1+O(\varphi^{-2m})\bigr).
\]
取二进对数即得
\[
B_m
\ge
m\log_2\!\Bigl(\frac{2}{\varphi}\Bigr)
+\log_2\!\Bigl(\frac{\sqrt5}{\varphi^2}\Bigr)
+o(1)
=
m\log_2\!\Bigl(\frac{2}{\varphi}\Bigr)+O(1).
\]

当 \(B_m\equiv 0\) 时，
\[
p_m\le \mathrm{Succ}_m(0)\le \bar d_m^{-1}\to 0,
\]
因为 \(\bar d_m\asymp (2/\varphi)^m\to +\infty\)。故纯标量通道不可能承载远位局域提取。证毕。
\end{proof}

\leanverified{paper\_conclusion\_foldgauge\_pi\_sufficient\_statistic\_fiber\_obstruction}
\begin{theorem}[逆代换充分统计外置的纤维不可逆障碍]\label{thm:conclusion-foldgauge-pi-sufficient-statistic-fiber-obstruction}
固定 \(m\ge 1\) 与 \(x\in X_m\)。设
\[
\mathcal R_m:\Omega_m\longrightarrow \mathcal A_m
\]
为任意外置残差，并假定其在纤维 \(\Fold_m^{-1}(x)\) 上只依赖于逆代换次数 \(r(\omega)\)。记
\[
Z_x(t):=\sum_{\omega\in \Fold_m^{-1}(x)}t^{r(\omega)},
\qquad
d_m(x):=\bigl|\Fold_m^{-1}(x)\bigr|.
\]
则必有
\[
\bigl|\mathcal R_m(\Fold_m^{-1}(x))\bigr|
\le
\deg Z_x+1.
\]
因此，只要
\[
d_m(x)>\deg Z_x+1,
\]
映射
\[
\omega\longmapsto \bigl(\Fold_m(\omega),\mathcal R_m(\omega)\bigr)
\]
便不可能在该纤维上单射。进一步，在分量分解
\[
\Gamma_x\simeq\bigsqcup_i P_{\ell_i}
\]
下有闭式
\[
d_m(x)=\prod_i F_{\ell_i+2},
\qquad
\deg Z_x=\sum_i \left\lceil\frac{\ell_i}{2}\right\rceil,
\]
故除去有限退化类型外，任何把 \(\pi\)-方案压缩到“每纤维一个逆代换充分统计量”的实现，都不能同时保持精确性与纤维可逆性。
\end{theorem}
\begin{proof}
这正是命题 \ref{prop:pom-sufficient-statistic-residual-noninvertibility} 在 fold-\(\pi\) 语境下的直接专门化。命题已经给出像集大小上界
\[
\bigl|\mathcal R_m(\Fold_m^{-1}(x))\bigr|
\le
\deg Z_x+1,
\]
并证明当 \(d_m(x)>\deg Z_x+1\) 时，\(\omega\mapsto (\Fold_m(\omega),\mathcal R_m(\omega))\) 在该纤维上不可能单射。末尾的 Fibonacci 乘积闭式与次数公式也同样来自该命题。证毕。
\end{proof}

\begin{theorem}[AGM 型翻译器的局部型别与地址预算二歧]\label{thm:conclusion-foldgauge-pi-agm-translator-dichotomy}
\leanverified{paper\_conclusion\_foldgauge\_pi\_agm\_translator\_dichotomy}
设 \(\Sigma\) 为有限字母表，\(|\Sigma|=M\ge 2\)。则不存在一族把 AGM/Gauss--Salamin 型 \(\pi\) 通道嵌入 fold-\(\pi\) 读出的翻译器同时满足下列两项：
\begin{enumerate}
  \item 在边界固定点附近把 AGM 误差芽局部解析共轭到 fold-gauge \(\pi\) 通道的误差芽；
  \item 存在确定性可逆提升
  \[
  \Psi_m:\Omega_m\longrightarrow X_m\times \Sigma^{\le L_m},
  \qquad
  \mathrm{pr}_{X_m}\circ \Psi_m=\Fold_m,
  \]
  且
  \[
  \limsup_{m\to\infty}\frac{L_m}{m}
  <
  \frac{\log(2/\varphi)}{\log M}.
  \]
\end{enumerate}
等价地，任何 AGM--fold 翻译器若坚持确定性可逆提升，就必须满足线性密度级侧信息下界
\[
\liminf_{m\to\infty}\frac{L_m}{m}
\ge
\frac{\log(2/\varphi)}{\log M},
\]
而若试图把该代价压到更低，则只能放弃局部解析正规形的一致性。
\end{theorem}
\begin{proof}
由定理 \ref{thm:conclusion-foldgauge-pi-koenigs-vs-agm-nonconjugacy}，fold-gauge \(\pi\) 误差芽属于非零乘子的 Koenigs 型，而 AGM/Gauss--Salamin 误差芽属于乘子为 \(0\) 的 B\"ottcher 型，故二者不可能局部解析共轭。这排除了第一项。

另一方面，若存在确定性可逆提升 \(\Psi_m\)，则定理 \ref{thm:conclusion-side-info-length-lower-bound} 直接给出
\[
L_m\ge \left\lceil \log_M\!\Bigl(\frac{2^m}{F_{m+2}}\Bigr)\right\rceil-1,
\]
从而
\[
\liminf_{m\to\infty}\frac{L_m}{m}
\ge
\frac{\log(2/\varphi)}{\log M}.
\]
因此，只要要求 \(\limsup_m L_m/m<\log(2/\varphi)/\log M\)，可逆提升就不可能存在。两项合并即得结论。证毕。
\end{proof}

\begin{theorem}[零外置容量同一性蕴含 Bernoulli--偶 \texorpdfstring{$\zeta$}{zeta} 完备]\label{thm:conclusion-foldgauge-pi-zero-external-capacity-bernoulli-completeness}
设 \(F\) 为任意零外置的纯标量 \(\pi\) 通道，并记其在尺度 \(m\) 上诱导的连续容量曲线为
\[
\mathcal C_m^{F,\mathrm{cont}}(T).
\]
若对一切充分大的 \(m\) 与所有 \(T\ge 0\) 都有
\[
\mathcal C_m^{F,\mathrm{cont}}(T)
=
\mathcal C_m^{\mathrm{fold},\mathrm{cont}}(T),
\]
则对每个 \(q>0\)，其全部正阶矩谱满足
\[
S_q^F(m)=S_q^{\mathrm{fold}}(m).
\]
更进一步，同一条容量曲线序列唯一决定 fold-gauge 规范常数 \(C_m\) 的整条 Stirling--Bernoulli 渐近谱，因而唯一决定全部偶 Bernoulli 数 \(B_{2r}\) 与全部偶值 \(\zeta(2r)\)。特别地，任何不编码完整偶 \(\zeta\)-塔的纯标量 \(\pi\) 方案，都不可能属于 fold-\(\pi\) 的零外置容量等价类。
\end{theorem}
\begin{proof}
由定理 \ref{thm:conclusion-capacity-equivalence-thermodynamic-rigidity}，容量曲线逐点相同蕴含
\[
S_q^F(m)=S_q^{\mathrm{fold}}(m)
\qquad(q>0),
\]
故全部正阶矩谱一致。

另一方面，定理 \ref{thm:conclusion-capacity-curve-mellin-complete-reconstruction} 已证明：单独给定整条连续容量曲线 \(\mathcal C_m^{\mathrm{cont}}(\cdot)\)，便可唯一恢复全部负奇 Mellin 平均 \(A_{m,2r-1}\)，并因而唯一恢复 fold-gauge 规范常数列 \((C_m)_{m\ge 1}\) 的整个 Stirling--Bernoulli 渐近谱。再由定理 \ref{thm:fold-bin-gauge-constant-stirling-bernoulli-hierarchy} 与命题 \ref{prop:fold-bin-gauge-bernoulli-extraction-operator}，这些系数逐阶唯一决定全部 \(B_{2r}\)；结合
\[
B_{2r}=(-1)^{r-1}\frac{2(2r)!}{(2\pi)^{2r}}\zeta(2r),
\]
便得到全部偶值 \(\zeta(2r)\)。证毕。
\end{proof}

\begin{proposition}[BBP 型十六进制 lift 的最小地址率]\label{prop:conclusion-foldgauge-pi-bbp-hex-address-rate}
\leanverified{paper\_conclusion\_foldgauge\_pi\_bbp\_hex\_address\_rate}
若某个 BBP 型十六进制位抽取方案被实现为 fold 语义下的确定性可逆提升，并且每一步只写入一个十六进制符号，即
\[
|\Sigma|=16,
\]
则其侧信息长度必满足
\[
\liminf_{m\to\infty}\frac{L_m}{m}
\ge
\frac{\log(2/\varphi)}{\log 16}
=
0.07643952159234564\ldots\,.
\]
\end{proposition}
\begin{proof}
对定理 \ref{thm:conclusion-foldgauge-pi-agm-translator-dichotomy} 中的线性密度下界取 \(M=16\) 即得。等价地，也可直接对定理 \ref{thm:conclusion-side-info-length-lower-bound} 代入 \(|\Sigma|=16\)。证毕。
\end{proof}

\endinput

\begin{theorem}[折叠双谱对偶决定偶部 Archimedean 塔]\label{thm:conclusion-foldgauge-double-spectrum-archimedean-tower}
\leanverified{paper\_conclusion\_foldgauge\_double\_spectrum\_archimedean\_tower}
设 \(\{r_q\}_{q\ge 2}\) 为碰撞 moment-kernel 的经审计 Perron 根序列，而 \(\{C_m\}_{m\ge 1}\) 为 fold-gauge 规范常数序列。则二者唯一决定
\[
\varphi,\qquad
2\pi,\qquad
\{B_{2r}\}_{r\ge 1},\qquad
\{\zeta(2r)\}_{r\ge 1}.
\]
更精确地，
\[
\varphi=\left(\lim_{q\to\infty}r_q^{1/q}\right)^2,
\qquad
2\pi=\lim_{m\to\infty}C_m.
\]
若进一步记
\[
a_r:=
\frac{B_{2r}}{r(2r-1)}
\bigl(\varphi^{-1}+\varphi^{2r-3}\bigr),
\]
则对每个整数 \(r\ge 1\) 都有递归恢复公式
\[
a_r=
\lim_{m\to\infty}
\left(\frac{2}{\varphi}\right)^{(2r-1)m}
\left(
\log C_m-\log(2\pi)-\sum_{j=1}^{r-1}a_j\left(\frac{\varphi}{2}\right)^{(2j-1)m}
\right),
\]
从而
\[
B_{2r}
=
\frac{r(2r-1)}{\varphi^{-1}+\varphi^{2r-3}}\,a_r.
\]
再由
\[
B_{2r}=(-1)^{r-1}\frac{2(2r)!}{(2\pi)^{2r}}\zeta(2r),
\]
便可逐阶恢复全部 \(\zeta(2r)\)。
\end{theorem}
\begin{proof}
由命题 \ref{prop:pom-rq-qinfty-endpoint}，
\[
\lim_{q\to\infty}r_q^{1/q}=\sqrt{\varphi},
\]
故 \(\{r_q\}_{q\ge 2}\) 唯一决定 \(\varphi\)。另一方面，推论 \ref{cor:fold-bin-recover-2pi} 给出
\[
C_m\to 2\pi,
\]
故 \(\{C_m\}_{m\ge 1}\) 唯一决定 \(2\pi\)。

再由定理 \ref{thm:conclusion-bernoulli-zeta-tower-explicit-triangular-inversion}，一旦 \(\varphi\) 与 \(2\pi\) 已知，\((\log C_m)\) 的 Stirling--Bernoulli 层级便按上述极限公式逐阶唯一恢复全部 \(a_r\)，从而唯一恢复全部 \(B_{2r}\)；最后代入 Bernoulli--\(\zeta\) 恒等式即得全部 \(\zeta(2r)\)。证毕。
\end{proof}

\begin{corollary}[两尺度谱刚性]\label{cor:conclusion-foldgauge-two-scale-spectral-rigidity}
若两套折叠系统具有完全相同的 Perron 根序列 \(\{r_q\}_{q\ge 2}\) 与规范常数序列 \(\{C_m\}_{m\ge 1}\)，则它们诱导的 Stirling--Bernoulli 形式射流完全一致。等价地，它们的形式 Archimedean 因子
\[
\log \Gamma(z)\sim
\left(z-\frac12\right)\log z-z+\frac12\log(2\pi)
+\sum_{r\ge 1}\frac{B_{2r}}{2r(2r-1)}\,z^{-(2r-1)}
\]
在无穷远处逐项相同。因而，任何同时保持这两条谱轴的变形，都不能改变任何偶值 \(\zeta(2r)\)。
\end{corollary}
\begin{proof}
由定理 \ref{thm:conclusion-foldgauge-double-spectrum-archimedean-tower}，\(\{r_q\}_{q\ge 2}\) 与 \(\{C_m\}_{m\ge 1}\) 联合唯一决定
\[
\varphi,\qquad
2\pi,\qquad
\{B_{2r}\}_{r\ge 1},\qquad
\{\zeta(2r)\}_{r\ge 1}.
\]
而 Stirling 形式射流只依赖于 \(2\pi\) 与全部 \(B_{2r}\)，故结论成立。证毕。
\end{proof}

\[
V_6(w):=\sum_{k=1}^{6}w_kF_{k+1},
\qquad
\beta(a,b):=21a+34b,
\]
\[
I_2=\{(0,0),(0,1)\},
\qquad
I_3=\{(0,0),(1,0),(0,1)\},
\qquad
I_4=\{0,1\}^2.
\]

\begin{theorem}[window-\texorpdfstring{$6$}{6} dyadic 区间的四层 hidden-strip 精确剖分]\label{thm:conclusion-window6-hidden-strip-fourlayer-partition}
\leanverified{paper\_conclusion\_window6\_hidden\_strip\_fourlayer\_partition}
定义闭区间
\[
J_{00}:=[0,20],
\qquad
J_{10}:=[21,33],
\qquad
J_{01}:=[34,54],
\qquad
J_{11}:=[55,63].
\]
则有严格分解
\[
\{0,1,\dots,63\}=J_{00}\sqcup J_{10}\sqcup J_{01}\sqcup J_{11},
\]
并且更强地，
\[
J_{00}=\{V_6(w):w\in X_6\},
\]
\[
J_{10}=\{V_6(w)+21:\ d_6^{\mathrm{bin}}(w)\ge 3\},
\]
\[
J_{01}=\{V_6(w)+34:\ w\in X_6\},
\]
\[
J_{11}=\{V_6(w)+55:\ d_6^{\mathrm{bin}}(w)=4\}.
\]
等价地，
\[
V_6(S_{6,4})=\{0,\dots,8\},
\qquad
V_6(S_{6,3})=\{9,10,11,12\},
\qquad
V_6(S_{6,2})=\{13,\dots,20\}.
\]
\end{theorem}
\begin{proof}
文稿已给出 \(V_6:X_6\to\{0,1,\dots,20\}\) 为双射，故
\[
V_6(X_6)=\{0,1,\dots,20\}.
\]
又由 rigid \(X_4\) 子覆盖描述，
\[
S_{6,2}=\{(v,0,1):v\in X_4\},
\]
从而
\[
V_6(S_{6,2})=13+V_4(X_4)=13+\{0,\dots,7\}=\{13,\dots,20\}.
\]
另一方面，window-\(6\) 的四个三点纤维对应于
\[
\texttt{001010},\qquad
\texttt{010010},\qquad
\texttt{100010},\qquad
\texttt{101010},
\]
代入 Fibonacci 权 \((1,2,3,5,8,13)\) 可得
\[
V_6(\texttt{001010})=11,\qquad
V_6(\texttt{010010})=10,\qquad
V_6(\texttt{100010})=9,\qquad
V_6(\texttt{101010})=12.
\]
故
\[
V_6(S_{6,3})=\{9,10,11,12\}.
\]
于是
\[
V_6(S_{6,4})
=
\{0,\dots,20\}\setminus\bigl(V_6(S_{6,2})\cup V_6(S_{6,3})\bigr)
=
\{0,\dots,8\}.
\]
再由文稿给出的仿射 Boolean 方块标准形，
\[
(\Fold_6^{\mathrm{bin}})^{-1}(w)=V_6(w)+\beta(I_{d_6^{\mathrm{bin}}(w)}),
\]
并且
\[
\beta(I_2)=\{0,34\},
\qquad
\beta(I_3)=\{0,21,34\},
\qquad
\beta(I_4)=\{0,21,34,55\}.
\]
因此 \(0\) 与 \(34\) 两个偏移对所有纤维同时出现，\(21\) 只在 \(d_6^{\mathrm{bin}}(w)\ge 3\) 时出现，而 \(55\) 只在 \(d_6^{\mathrm{bin}}(w)=4\) 时出现。结合
\[
V_6(S_{6,3}\cup S_{6,4})=\{0,\dots,12\},
\qquad
V_6(S_{6,4})=\{0,\dots,8\},
\]
便得到四个 strip 的精确描述。四段区间长度分别为 \(21,13,21,9\)，总和为 \(64\)，故给出 \(\{0,\dots,63\}\) 的不交剖分。证毕。
\end{proof}

\begin{corollary}[window-\texorpdfstring{$6$}{6} 隐方块对 dyadic 几何级数的望远镜分解]\label{cor:conclusion-window6-dyadic-geometric-series-hidden-strip-telescope}
\leanverified{paper\_conclusion\_window6\_dyadic\_geometric\_series\_hidden\_strip\_telescope}
在形式幂级数环中有恒等式
\[
\sum_{n=0}^{63}z^n
=
\sum_{n=0}^{20}z^n
+\sum_{n=21}^{33}z^n
+\sum_{n=34}^{54}z^n
+\sum_{n=55}^{63}z^n,
\]
并且可重写为
\[
\frac{1-z^{64}}{1-z}
=
(1+z^{34})\frac{1-z^{21}}{1-z}
+z^{21}\frac{1-z^{13}}{1-z}
+z^{55}\frac{1-z^{9}}{1-z}.
\]
\end{corollary}
\begin{proof}
第一式只是定理 \ref{thm:conclusion-window6-hidden-strip-fourlayer-partition} 的区间分解在几何级数上的逐项改写。
对第二式，只需将右端乘以 \(1-z\)：
\[
(1-z^{21})+(z^{34}-z^{55})+(z^{21}-z^{34})+(z^{55}-z^{64})
=
1-z^{64}.
\]
故结论成立。证毕。
\end{proof}

\begin{theorem}[window-\texorpdfstring{$6$}{6} 上 \texorpdfstring{$e$}{e} 的前 \texorpdfstring{$64$}{64} 项有限 Borel 剖分]\label{thm:conclusion-window6-e64-finite-borel-partition}
\leanverified{paper\_conclusion\_window6\_e64\_finite\_borel\_partition}
对每个 \(t\in\CC\)，定义
\[
E_{64}(t):=\sum_{n=0}^{63}\frac{t^n}{n!}.
\]
则有精确恒等式
\[
E_{64}(t)
=
\sum_{v=0}^{8}\ \sum_{u\in I_4}\frac{t^{v+\beta(u)}}{(v+\beta(u))!}
+\sum_{v=9}^{12}\ \sum_{u\in I_3}\frac{t^{v+\beta(u)}}{(v+\beta(u))!}
+\sum_{v=13}^{20}\ \sum_{u\in I_2}\frac{t^{v+\beta(u)}}{(v+\beta(u))!}.
\]
等价地，
\[
E_{64}(t)
=
\sum_{n=0}^{20}\frac{t^n}{n!}
+\sum_{n=21}^{33}\frac{t^n}{n!}
+\sum_{n=34}^{54}\frac{t^n}{n!}
+\sum_{n=55}^{63}\frac{t^n}{n!}.
\]
特别地，
\[
E_{64}(1)
=
\sum_{v=0}^{8}
\left(
\frac1{v!}+\frac1{(v+21)!}+\frac1{(v+34)!}+\frac1{(v+55)!}
\right)
\]
\[
\phantom{E_{64}(1)=}
+\sum_{v=9}^{12}
\left(
\frac1{v!}+\frac1{(v+21)!}+\frac1{(v+34)!}
\right)
+\sum_{v=13}^{20}
\left(
\frac1{v!}+\frac1{(v+34)!}
\right),
\]
并满足
\[
0<e-E_{64}(1)<\frac{2}{64!}.
\]
\end{theorem}
\begin{proof}
由定理 \ref{thm:conclusion-window6-hidden-strip-fourlayer-partition}，\(\{0,\dots,63\}\) 被精确分成四层 hidden strips；将数列
\[
a_n:=\frac{t^n}{n!}
\]
沿该剖分逐项求和，立得第二个显示公式。再用
\[
\{0,\dots,8\}\times I_4,\qquad
\{9,\dots,12\}\times I_3,\qquad
\{13,\dots,20\}\times I_2
\]
对四层 strip 做原像重组，即得第一个显示公式。

当 \(t=1\) 时，第三个显示式只是把上述分解写成显式阶乘和。另一方面，
\[
e=\sum_{n=0}^{\infty}\frac1{n!},
\]
故
\[
e-E_{64}(1)=\sum_{n=64}^{\infty}\frac1{n!}
\le
\frac1{64!}\sum_{j=0}^{\infty}\frac1{65^j}
<
\frac{2}{64!}.
\]
显然该余项严格为正。证毕。
\end{proof}


\begin{conjecture}[fold-\texorpdfstring{$\pi$}{pi} 普适类的三不变量判别]\label{conj:conclusion-foldgauge-modular-acceleration-measurability}
记
\[
\mathscr F_C:=\sigma(C_m:\ m\ge 1)
\]
为由折叠标量序列生成的可见标量代数。设一个投影兼容的 \(\pi\)-通道在每个尺度上都允许纯标量 fold 读出。则其与 fold-\(\pi\) 落在同一普适类，当且仅当同时满足下列三项：
\begin{enumerate}
  \item 奇层可消元性：其有限窗加速器可按定理 \ref{thm:conclusion-foldgauge-pi-oddlayer-annihilation-hierarchy} 的方式逐层湮灭奇次污染模态；
  \item Bernoulli 完备性：其零外置容量类在定理 \ref{thm:conclusion-foldgauge-pi-zero-external-capacity-bernoulli-completeness} 的意义下唯一恢复全部 \(B_{2r}\) 与全部 \(\zeta(2r)\)；
  \item 边界乘子非零：其边界误差芽落在定理 \ref{thm:conclusion-foldgauge-pi-koenigs-vs-agm-nonconjugacy} 所刻画的非零乘子 Koenigs 型，而不进入 AGM 式的超吸引平方化类。
\end{enumerate}
预期地，纯标量 Machin-like lift 将被定理 \ref{thm:conclusion-foldgauge-pi-sufficient-statistic-fiber-obstruction} 所揭示的纤维不可逆障碍排除；AGM 类将被定理 \ref{thm:conclusion-foldgauge-pi-agm-translator-dichotomy} 的局部型别--地址预算分裂排除；BBP 类将被命题 \ref{prop:conclusion-foldgauge-pi-bbp-hex-address-rate} 的正地址率墙排除；而 Ramanujan--Chudnovsky 类若要进入该普适类，则必须携带某个真正超出 \(\mathscr F_C\) 的模形式辅助模态，而不能停留在纯标量投影之内。
\end{conjecture}

\paragraph{比较读法}
上述命题链说明：折叠标量通道并不只把 \(2\pi\) 作为极限值锁定，而是与碰撞 Perron 谱轴共同组成一条双谱刚性接口：\(\{r_q\}_{q\ge 2}\) 锁定 \(\varphi\)，\(\{C_m\}_{m\ge 1}\) 再把偶 Bernoulli--\(\zeta\) 塔压缩为同一条奇次黄金幂谱链。于是，纯标量层上的 \(\pi\)-逼近速率被定理 \ref{thm:conclusion-foldgauge-pi-golden-linear-law} 刚性固定；不同投影兼容序列的首个可观测分歧，由定理 \ref{thm:conclusion-foldgauge-pi-evenzeta-first-difference} 与偶 \(\zeta\)-塔的首差阶严格对位；而一旦把可逆提升、容量同一与外部经典机制并读，则进一步出现一条更强的分类链：定理 \ref{thm:conclusion-foldgauge-pi-sufficient-statistic-fiber-obstruction} 表明任何把每纤维残差压缩成单一逆代换充分统计量的纯标量 lift 都会触发结构性不可逆；定理 \ref{thm:conclusion-foldgauge-pi-agm-translator-dichotomy} 说明 AGM 类若欲进入 fold 语义，必须在局部正规形保持与线性密度侧信息之间二择其一；定理 \ref{thm:conclusion-foldgauge-pi-zero-external-capacity-bernoulli-completeness} 与定理 \ref{thm:conclusion-foldgauge-double-spectrum-archimedean-tower} 则表明零外置容量同一性连同碰撞 Perron 端点已共同强制携带整条 Bernoulli--偶 \(\zeta\) 谱塔；命题 \ref{prop:conclusion-foldgauge-pi-bbp-hex-address-rate} 进一步把 BBP 的十六进制局域性转译为不可消除的正地址率。就此而言，本文内部所谓 \(\pi\) 通道并非孤立数值算法，而是一个以奇层消元、双谱 Bernoulli 完备性、Koenigs 边界型别与正地址率障碍共同刻画的折叠标量普适类。
同一标量接口在 window-\(6\) 上还把 dyadic 区间 \(\{0,\dots,63\}\) 精确剖分为四层 hidden strips，并由此把 \(e\) 的前 \(64\) 项截断压缩成一个有限 Borel 分块；于是 cubical hidden register 与标量指数级数之间也出现了与 Bernoulli--\(\zeta\) 链平行的另一条刚性对位。

\endinput
